\chapter{Inter-core communication}\label{ch:intercore}

Cores share no writable memory. They cooperate by passing capabilities: a core that wants another core
to see data hands over the capability for it. The hand-over is both the communication and the
synchronisation.

\section{Nodes and inboxes}\label{sec:inbox}

Nodes $0 \ldots N-1$ are the cores, where $N$ is the number of cores, and node $N$ is block device~0
(Section~\ref{sec:blkdev}). Every core has an \term{inbox} that keeps a separate FIFO queue, two entries
deep, for each node that may give to it. The block device has a single four-entry inbox.

\begin{description}[leftmargin=1em,style=nextline]
\item[\insn{cgive cs, rx}] appends \reg{cs} to the giver's own queue in the inbox of node \reg{rx}. A
  linear capability leaves \reg{cs} empty; a copyable one is copied. If the giver's queue at that node is
  full, the giver waits; gifts from other senders never make it wait. Giving to a core that has halted,
  or to a node number above $N$, traps with \cause{capop}.
\item[\insn{ctake cd, rs}] removes the oldest gift \emph{from node} \reg{rs} into \reg{cd}, waiting until
  there is one. Supervisor mode only, like the next form, \insn{cpoll} and the inbox bit of \insn{tany}
  (Section~\ref{sec:userinbox}).
\item[\insn{ctake cd}] removes the oldest gift of the \emph{lowest-numbered} sender that has one,
  waiting while the inbox is empty.
\item[\insn{csend cs, ce}, \insn{crecv cd, ce}] give to and take from the node named by an endpoint
  capability (Section~\ref{sec:endpoints}).
\item[\insn{cpoll rd}] returns a bit mask with bit $i$ set if a gift from sender $i$ is waiting. It
  never waits.
\item[\insn{tany rd, rs}] with bit~16 of \reg{rs} set waits until any gift is waiting, possibly together
  with transfer tags (Section~\ref{sec:events}).
\item[\csr{takefrom}] holds the sender of the capability most recently taken by \insn{ctake} or
  \insn{crecv}.
\end{description}

\section{Ordering guarantees}\label{sec:publication}

\begin{rule*}{The publication rule}
\insn{cgive}, \insn{csend}, \insn{cshare}, \insn{catomic} and \insn{tamo} take effect only when none of
the issuing core's memory transfers (\insn{tget}, \insn{tput} or \insn{tgetx} to T2 or T3) is in
flight. When the receiver's \insn{ctake} or \insn{crecv} returns, every byte the giver wrote with
\insn{tput}, through that capability or any other, is in far memory, and every byte the giver read with
\insn{tget} or \insn{tgetx} has been read. \insn{halt} waits for all of the core's transfers, device
transfers and pending device reads included, in user mode as in supervisor mode.
\end{rule*}

So a producer can fill a buffer with \insn{tput}s and immediately \insn{cgive} it; the consumer may
\insn{tget} from it as soon as it has taken it. No fences, flags or locks are involved, and since there
are no caches there is no coherence traffic. The \texttt{pipeline} and \texttt{race} examples rely on
this rule.

Device transfers do not hold a hand-over up. A device read can stay in flight indefinitely (a
pending console break read completes only on Ctrl-C), and device space is not memory that a
capability hand-over publishes, so a core may give, share or convert capabilities while it waits for a
device event. \insn{halt} is different: it drains everything, so a core that is about to halt first
retires a break read by writing the console's release port (Section~\ref{sec:break}). So is a gift to
the block device's node: it also waits for the core's device writes, so that the command it submitted
just before is at the device when its buffer arrives (Section~\ref{sec:blkdev}). Core-control writes
count as memory transfers here; besides, the next instruction waits for them (Section~\ref{sec:corectl}).

Gifts from one sender are taken in the order they were given. Gifts from different senders are
independent queues.

\section{Deterministic fan-in}

When several cores may give to the same core at once, an any-sender \insn{ctake} returns the gift of
the lowest-numbered sender \emph{that has already delivered}, which depends on timing. To collect
results deterministically, take from each sender explicitly:

\begin{asm}
        li    s1, 1                     ; first worker
collect:
        ctake c5, s1                    ; only a gift from core s1 counts
        cjoin c4, c5                    ; append its piece (pieces arrive in core order)
        addi  s1, s1, 1
        blt   s1, s2, collect           ; s2 = number of cores
\end{asm}

Because every sender has its own queue, gifts that other senders deliver early can never fill the
inbox and block the sender a \insn{ctake cd, rs} is waiting for. Each (sender, receiver) pair behaves as
an independent bounded FIFO channel, which is what the determinism argument of
Section~\ref{sec:determinism} needs.

\section{Supervised execution}\label{sec:supervised}

An operating system runs a program written for the bare machine, which assumes it is core 0 of $N$ and
that \insn{halt} ends it, on a subset of the cores, in user mode. The core window and the privilege
modes make this possible without changing the program:

\begin{description}[leftmargin=1em,style=nextline]
\item[\csr{corebase}, \csr{corecount}] define a window of physical cores
  $[\textit{corebase}, \textit{corebase} + \textit{corecount})$. When \csr{corecount} is nonzero:
  \begin{itemize}
  \item \csr{coreid} reads as physical id $-$ \csr{corebase}, and \csr{ncores} reads as \csr{corecount};
  \item \insn{cgive} target \reg{rx} must be less than \csr{corecount} (else \cause{capop}) and is
        delivered to physical node \reg{rx} + \csr{corebase};
  \item \insn{ctake cd, rs} matches sender \reg{rs} + \csr{corebase}, and \csr{takefrom} and
        \insn{cpoll} report senders relative to \csr{corebase} (these two instructions execute only in
        supervisor mode, where the supervisor may run them on a program's behalf).
  \end{itemize}
  \csr{corebase} applies even when \csr{corecount} is zero; only the window check and \csr{ncores} depend
  on \csr{corecount} being nonzero. Both CSRs are writable only in supervisor mode.
\item[User mode] The supervisor enters the program with \insn{eret} and \csr{status}.SPP = 0. The
  program's \insn{halt} traps back to the supervisor with cause \cause{halt}
  (Section~\ref{sec:halt}); its faults go to its own handler if it installed one in \csr{tvec} (which
  in user mode is \csr{utvec}), and otherwise to the supervisor.
\end{description}

A supervisor sets the window on each core of the program (in code it runs on that core before entering
the program), gives the program only the capabilities it should have, and points its own \csr{tvec} at
its handler. In user mode \insn{cgive} cannot leave the window and the window cannot be changed, so
the program cannot address cores outside it, nor the block device unless the window includes node $N$.
The capabilities it holds are its only authority over memory and devices; the capability save table,
the revocation operations and the creation of endpoints are out of its reach. To let the program reach a
node outside its window, the supervisor gives it an endpoint for that node. To take its memory back
regardless of where the program has passed it, the supervisor lends it revocable memory. To take its core back, the supervisor sends a stop request through core control
(Section~\ref{sec:corectl}).

\subsection{The inbox in user mode}\label{sec:userinbox}

\begin{rule*}{The inbox belongs to the supervisor}
In user mode \insn{ctake} (with or without a sender), \insn{cpoll}, and \insn{tany} with bit~16 of its
operand set always trap with cause \cause{priv}. \csr{tval} is the instruction word.
\end{rule*}

A core's inbox receives everything that is given to the core: gifts for the program, and whatever the
supervisor itself is sent, such as requests, replies and memory for the next program. If a program
could take from the inbox it could take what was meant for its supervisor, or for another program that
shares the core. So the supervisor takes every arrival and decides whose it is, and the three inbox
instructions trap. Three things remain available to a user program:
\begin{itemize}
\item \insn{cgive}, restricted to the core window, and \insn{csend} through an endpoint with \perm{W};
\item \insn{crecv} through an endpoint with \perm{R}, which names one sending node; a supervisor that
  keeps all receiving to itself simply never hands out \perm{R};
\item \insn{tany}, \insn{twait} and \insn{tpoll} on transfer tags.
\end{itemize}

The \cause{priv} trap is also what lets a supervisor run programs written for the bare machine
unchanged: it \emph{emulates} the trapping instruction. The handler decodes the instruction word in
\csr{tval}, performs the take or poll on the queue of gifts it keeps for that program, writes the
result to the program's registers and returns to \csr{epc} + 4; if nothing has arrived, it suspends the
program, runs another, and retries the instruction later. StrataOS does exactly this
(Section~\ref{sec:os-emul}), which is how the multicore examples run as time-shared threads.
